
\section{Realization} 
{\large
	The first part of this chapter provides an overview of the information flow between each entity, and the tasks of the given entitys. Aswell as an close look into the implementation of the functional requirements, by analysing code snippets, aswell as explaining the usage of the final product. While the second part will explain every necessity to assemble all parts together for a working product.  
	
}
\subsection{Definiton of Entitys} \label{entityDef}
{\large  As shown in table \ref{tab:entitysDefinition} each entity has its scope of dutys, making the final code more modular and consequently maintainable.}
\begin{table*}[!htpb]
	\begin{center}
		\caption{Entiys area of responsibility}
		\label{tab:entitysDefinition}
		\begin{tabular}{l|l} % <-- Alignments: 1st column left, 2nd middle and 3rd right, with vertical lines in between
			\textbf{Layer} & \textbf{Functionality} \\
			\hline & \\[-1.5ex]
			BLEViewModel & \makecell[l]{
				- Communication with the microcontroller\\-  Receives datastream\\-  Sends instructions to microcontroller\\
			}\\
			\hline & \\[-1.5ex]
			GraphEcgWidget & \makecell[l]{Draws the ECG data in realtime,\\ and displays Beats per minute}  \\
			\hline& \\[-1.5ex]
			Storewrapper& \makecell[l]{Converts data from microcontroller\\ into drawble data } \\
			\hline& \\[-1.5ex]
			BPM& \makecell[l]{Calculates the current Beats Per Minute}\\
			\hline& \\[-1.5ex]
			VitalMonitorWidet & \makecell[l]{Plots up to 24 Hour long  recoreded ECG \\data and the BPM of the plotted Intervall. \\
				Provides zooming functionality.} \\
			\hline& \\[-1.5ex]
			StorageHandler & \makecell[l]{Allows usage of the OS filechooser,\\ to select files.}\\
			\hline& \\[-1.5ex]
			MainView & \makecell[l]{Collection of visble elements (Widgets),\\ to provide reusability}\\
			\hline & \\[-1.5ex]
			MainActivity& \makecell[l]{
				Visible to user, contains the complete\\ User Interface. Takes and processes user\\ inputs.
			}\\
			\hline
			\end{tabular}
		\end{center}
	\end{table*}
	\newpage
{\large
	

}

\subsection{Inter-Entity Communication Overview }
{\large
	The following diagramm shows the Information flow between all entitys. The unidirectional arrow means that, information is flowing oneside. An informationexchange between two entitys is represented by an bidirectional arrow.
	
	\begin{center}
	\includegraphics[scale = 0.70]{/Users/eric/Downloads/InformationFlow.drawio-3.png}
	\label{fig:informationFlowDiag}
	\captionof{figure}{Informationflow Diagramm}
	\end{center}
	
	As already stated in \ref{entityDef} the BLEViewModel entity, serves as an interface between controller and app. Therfore every entity that needs to evaluate recorded data, is dependent on this Instance. Two of the entities listed in \ref{entityDef} are affected. To be more precise, it is the StoreWrapper that makes the data drawable (1) and forwards the data to the GraphECGWidget (4).
	The BPM instance that calculates the beat frequency (2), (which is  then displayed by the GraphECGWidget (5)).
	When comparing Informationflow (4) and (6) both seem very similiar, because both widgets obtain drawable data. However the Informationflow of (6) is bidirectional.  This is necessary because  the recorded data needs to be read, by the StorageHandler from a local file (7), before being sent to the Storewrapper (6). While in (4) the Storewrapper receives data directly from the ECG. After the data has been mapped by the Storewrapper, it is sent back to the VitalMonitorWidget (6).
	MainView integrates both the GraphECGWidget (8) and the VitalMonitorWidget (9) as components and manages their visibility.
	However, the range of tasks of the BLEViewModel is not limited to exchanging data, but also extends to scanning and connecting with other devices in the vicinity, as well as integrating and managing a 24-hour ECG recording. These functionalities are utilized by the MainView and the MainActivity. As already mentioned in table \ref{tab:entitysDefinition}, MainView serves solely to construct components of functionalities that are later utilized by MainActivity, such as buttons enabling users to start a 24-hour ECG recording (10). The MainActivity entity acts as a bridge between the microcontroller and the user, allowing device scanning, initiating recordings, establishing connections, and displaying ECG information (11).
	
\paragraph{Information Flow and Responsibilities of BLEViewModel}

\begin{itemize}
	\item \textbf{BLEViewModel as an Interface}  
	\begin{itemize}
		\item Acts as an interface between the controller and the app (\ref{entityDef}).  
		\item Any entity that evaluates recorded data depends on this instance.  
	\end{itemize}
	
	\item \textbf{Entities Dependent on BLEViewModel}  
	\begin{itemize}
		\item \textbf{StoreWrapper}  
		\begin{itemize}
			\item Makes the data drawable (1).  
			\item Forwards the data to the \textbf{GraphECGWidget} (4).  
		\end{itemize}
		\item \textbf{BPM Instance}  
		\begin{itemize}
			\item Calculates the beat frequency (2).  
			\item The result is displayed by the \textbf{GraphECGWidget} (5).  
		\end{itemize}
	\end{itemize}
	
	\item \textbf{Comparison of Information Flows (4) and (6)}  
	\begin{itemize}
		\item Both provide drawable data to widgets.  
		\item \textbf{Difference}:  
		\begin{itemize}
			\item \textbf{(4) One-way flow} $\rightarrow$ StoreWrapper receives data directly from the ECG.  
			\item \textbf{(6) Bidirectional flow} $\rightarrow$ Required for reading recorded data from a local file (7) using the \textbf{StorageHandler} before sending it to the StoreWrapper.  
		\end{itemize}
		\item Once mapped by the StoreWrapper, the data is sent back to the \textbf{VitalMonitorWidget} (6).  
	\end{itemize}
	
	\item \textbf{MainView Responsibilities}  
	\begin{itemize}
		\item Integrates \textbf{GraphECGWidget} (8) and \textbf{VitalMonitorWidget} (9).  
		\item Manages the visibility of these components.  
	\end{itemize}
	
	\item \textbf{Extended BLEViewModel Responsibilities}  
	\begin{itemize}
		\item Beyond data exchange, BLEViewModel also:  
		\begin{itemize}
			\item Scans and connects to other nearby devices.  
			\item Integrates and manages 24-hour ECG recordings.  
		\end{itemize}
		\item These functionalities are used by \textbf{MainView} and \textbf{MainActivity}.  
	\end{itemize}
	
	\item \textbf{MainView vs. MainActivity}  
	\begin{itemize}
		\item \textbf{MainView}  
		\begin{itemize}
			\item Constructs components and functionalities for \textbf{MainActivity}.  
			\item Provides UI elements, such as buttons for starting a 24-hour ECG recording (10).  
		\end{itemize}
		\item \textbf{MainActivity}  
		\begin{itemize}
			\item Acts as a bridge between the \textbf{microcontroller} and the \textbf{user}.  
			\item Handles device scanning, recording initiation, connection establishment, and ECG data display (11).  
		\end{itemize}
	\end{itemize}
\end{itemize}

	
}



\subsection{BLEViewModel }
{\large 
 The following subsections explore the logical implementation and functionality of individual entities, using code snippets and application screenshots for illustration. Each subsection is divided into a section covering the implementation and another focusing on the application's usage, with both sections referring to corresponding screenshots. Framework specific code fragments will be explaing in the course of relevant section.	

Before presenting the implementation of the BLE component, it is essential to establish a fundamental theoretical understanding of BLE.
This section offers a comprehensive explanation of the Bluetooth Low Energy, shedding light on its core principles, while exploring the architecture of the BLE stack. 
}

\subsubsection{BLE Stack}
{\large
	The BLE protocol stack consists of three main layers: the Controller, the Host, and the Application. The Controller encompasses the Physical Layer and the Link Layer, while the Host manages higher-level functionalities such as the Logical Link Control and Adaptation Protocol (L2CAP), the Attribute Protocol (ATT), the Generic Attribute Profile (GATT), the Security Manager (SM), and the Generic Access Profile (GAP). Communication between the Host and the Controller follows the standardized Host Controller Interface (HCI). At the top of the stack, the Application layer enables program execution. \cite{BLEStack}
	
	\begin{center}
		\includegraphics[scale = 0.90]{/Users/eric/Downloads/bluetooth_le_stack.png}
		\label{fig:BLEStack}
		\captionof{figure}{BLE protocol stack (https://de.mathworks.com/help/bluetooth/gs/bluetooth-protocol-stack.html)}
	\end{center}
	
	\paragraph{Application Layer}\mbox{}\\
	The application layer serves as the direct user interface, defining profiles that enable interoperability between various applications.
	
	\paragraph{Host Layer}\mbox{}\\
	Includes the host-sided HCI aswell as L2CAP, ATT, GATT, SMP and GAP. 
	
	\paragraph{Generic Acces Profile (GAP)}\mbox{}\\
	Manages the connection setup aswell as security, interfaces directly with profiles and services of the Application Layer.
	
	\paragraph{Security Manager Protocol (SMP)} \mbox{}\\
	Provides multiple security algorithms for encrypting and decrypting data packets, making the connection secure.
	
	\paragraph{Generic Attribute Profile (GATT)}\mbox{}\\
	Defines services and characteristics provided by the GATT Server. The GATT Server in this case is the microcontroller, which holds the definitions of services and characteristics. These can be requested by the GATT Client (in this case the application).
	Services organize data into logical entities and contain specific data units known as characteristics. Each service can include one or more characteristics and is uniquely identified by a numeric UUID. This UUID can be either 16-bit for officially adopted BLE services or 128-bit for custom services. Whereas the Characteristic is the lowest level concept in GATT transactions, which defines data. Services aswell as Characterics are main interactions points to work with when working with the BLE periphery, thus it is important to understand these terminologies.
	
	\paragraph{ Attribute Protocol (ATT)}\mbox{}\\
	The ATT enables the exchange of attributes (characteristics). Attributes are are composed of four components, the type (128-bit UUID), an handle (16-bit indentifier) to reference attributes, rights of the attribute(reading, writing), and the value of the attribute. The following figure gives a good overview of the structure of the attributes.
	
	\begin{center}
	\includegraphics[scale = 0.90]{/Users/eric/Downloads/ble_att_representation.png}
	\label{fig:ATTstructure}
	\captionof{figure}{Attribute Structure (https://de.mathworks.com/help/bluetooth/gs/bluetooth-protocol-stack.html)}
	\end{center}
	
	\paragraph{Logical Link Control and Adaption Layer Protocol (L2CAP)}\mbox{}\\
	Works as an interface betwenn the HCI and the upper modules of the Host layer, by multiplexing data packets to the corresponding protocols. Also fragments and reassembles large packets.
	
	\paragraph{Host-Side Host-Controller-Interface (HCI) }\mbox{}\\
	Generally said the HCI is the inteface between the host and the controller. The HCI establishes commands and events for packet data transmission and reception. When sending data, it processes raw information into packets, enabling communication between the host and the controller.
	
	\paragraph{Controller-Side Host-Controller-Interface (HCI) }\mbox{}\\
	As mentioned in the last paragraph, the HCI is the interface between the host and the controller. However, analogous to the host side, the data received from the controller is extracted and forwarded to the host.
	
	\paragraph{Controller Layer}\mbox{}\\
	Includes the Bluetooth Le PHY, LL, and controller-side HCI
	
	\paragraph{Link Layer (LL)}\mbox{}\\
	The LL functions similarly to the medium access control (MAC) layer in the OSI model. In Bluetooth, it directly interacts with the Bluetooth LE PHY and controls the radio’s link state, determining whether a device operates as a Central, Peripheral, Advertiser, or Scanner.
	
	\paragraph{Bluetooth LE PHY}\mbox{}\\
	 The physical layer (BLE PHY) is a 1 Mbps GFSK-based radio with adaptive frequency hopping, operating in the unlicensed 2.4 GHz ISM (industrial, scientific, and medical) band.
}	
\subsubsection{Implementation}
\paragraph{Dependency versions and Library}\mbox{}\\
{\large Before delving into the implementation, it is essential to note that a library provided by Nordic Semiconductor was employed in the development process.\footnote{\url{https://github.com/NordicSemiconductor/Android-BLE-Library}} The utilized dependency version are shown in the figure below:
	\begin{center}
	\includegraphics[scale = 0.90]{/Users/eric/Downloads/NordicSemiconLibVersions.png}
	\label{fig:NoSemiVersions}
	\captionof{figure}{Library Versions}
	\end{center} 

\paragraph{Scanning devices}\mbox{}\\
First of all we are exploring the \textit{scanning functionality}, to register nearby devices.
The following code snippet demonstrates the implementation of a BLE scanning method using Jetpack Compose. The \textit{startScan} function initializes the scanning process by setting the \textit{scanState} to \textit{SCANNING} and resetting the list of discovered BLE devices. The scanning operation is performed using \textit{BleScanner}, with results processed by \textit{BleScanResultAggregator}. The detected devices are filtered to include only those with the name \textit{"THI-EKG"} and stored in a MutableState. The scan runs asynchronously within the viewModelScope, and after a delay of 2000 milliseconds, it is terminated by canceling \textit{scanJob} and updating the \textit{scanState} to \textit{DONE}.
The snippet also contains two annotations above the function definition, which dont add any functionality.
Nevertheless, a brief explanation follows to avoid confusion and ensure clarification:


\textbf{\underline{@RequiresApi(Build.VERSION\textunderscore CODES.UPSIDE\textunderscore DOWN\textunderscore CAKE)}}\\  
This annotation ensures that the \textit{startScan} function is only executed on devices running Android API level \textbf{34 (Android 14, Upside Down Cake)} or higher. If the function is called on a lower API level, the compiler will issue a warning, and runtime execution may fail.  

 \textbf{\underline{@SuppressLint("MissingPermission")  }}\\
This suppresses lint warnings related to missing Bluetooth permissions (such as \textit{BLUETOOTH\textunderscore SCAN} or \textit{ACCESS\textunderscore FINE\textunderscore LOCATION}). Normally, scanning for BLE devices requires specific permissions, and Android Studio would generate a warning if they are not explicitly checked. However, this annotation tells the lint tool to ignore such warnings, assuming the necessary permissions are already handled elsewhere in the code. Thus the permissions are check in the Main Acitivity entitiy, this warning can be ignored, and will appear more often. 


\begin{figure}[H]
	\centering
	\makebox[0pt]{%
	\includegraphics[width=0.9\paperwidth]{/Users/eric/Downloads/ScanningCodeSnippet.png}}
	\label{fig:ScanningCodeSnipet}
	\captionof{figure}{Scanning functionality code }
\end{figure} 




}
\paragraph{Refreshing }
{\large
\begin{figure}[b]
	\centering
	\makebox[0pt]{%
		\includegraphics[width=0.9\paperwidth]{/Users/eric/Downloads/RefreshingBLE.png}}
	\label{fig:RefreshingCodeSnipet}
	\captionof{figure}{Refresh functionality code }
	\end{figure}
 }



